\chapter{Die Omega-Theorie und die Maxwell-Gravitation}

\section{Einleitung}

Die ART beschreibt Gravitation als Krümmung der Raumzeit.

Die Omega-Theorie beschreibt Gravitation als Rotation.

Dies ist kein bloßer Formalismus. Es ist eine tiefere Einsicht: Die Krümmung der ART ist ein emergenter Effekt – die fundamentale Größe ist die Rotation \( \vec{\Omega} \).

Aus der Wirbelstruktur der Weber-Gravitation folgt die Maxwell-Gravitation – und damit die natürliche Erklärung von Gravitationswellen mit Dispersion.

\section{Die Wirbelstruktur der Weber-Gravitation}

Die Weber-Gravitation enthält einen geschwindigkeitsabhängigen Term:

\[
\vec{F}_{\text{tan}} = -\frac{GMm}{r^2} \cdot \frac{2(\dot{r} \cdot \vec{v})}{c^2} \vec{v}
\]

Dieser Term ist nicht-zentral. Er besitzt eine nichtverschwindende Rotation (Wirbelstärke).

Die Wirbeldichte ist:

\[
\vec{\omega} = \nabla \times \vec{F}_{\text{WG}}
= \frac{6GMm}{c^2 r^3} (\hat{r} \cdot \vec{v})(\hat{r} \times \vec{v})
\]

Diese Wirbelstruktur ist der Schlüssel zu allem.

\section{Die Omega-Theorie}

Die Omega-Theorie basiert auf vier Axiomen:

\subsection{Axiom 1: Nicht-Lokalität und Instantaneität}

Die Gravitationswirkung ist nicht-lokal und instantan. Es gibt keine Ausbreitung mit endlicher Geschwindigkeit; die Wechselwirkung ist simultan.

\subsection{Axiom 2: Das Rotationsfeld}

Die Gravitation wird durch ein fundamentales, axiales Vektorfeld \( \vec{\Omega}(\vec{r}, t) \) beschrieben. Dieses Feld kodiert die lokale Rotationsdichte des Raumes.

\subsection{Axiom 3: Aufspaltung von Ursache und Wirkung}

Die Schwere (Ursache) und die Trägheit (Wirkung) sind fundamental verschieden. Die radiale Anziehung wird durch ein skalares Potential \( \Phi \) beschrieben, die tangentialen Trägheitskräfte durch das Rotationsfeld \( \vec{\Omega} \).

\subsection{Axiom 4: DSTT-Kompatibilität}

Die Theorie ist so konstruiert, dass aus ihr die Existenz einer Zustandsvariablen \( \beta(\vec{r}, t) \) folgt, die das lokale Verhältnis von tangentialer zu gesamter Trägheit angibt, und für die \( \dot{\beta} > 0 \) gilt.

\subsection{Die Feldgleichungen}

Die Felder werden nicht-lokal und instantan aus den Quellen erzeugt. Im Vakuum gehorchen sie formal den Maxwell-Gleichungen:

\[
\begin{aligned}
\nabla \cdot \vec{\Omega} &= 0 \\
\nabla \times \vec{\Omega} &= \frac{4\pi G}{c^2} \vec{j} + \frac{1}{c^2} \frac{\partial}{\partial t} (-\nabla \Phi) \\
\nabla \cdot (-\nabla \Phi) &= -4\pi G \rho \\
\nabla \times (-\nabla \Phi) &= -\frac{\partial \vec{\Omega}}{\partial t}
\end{aligned}
\]

\section{Die Maxwell-Gravitation}

Aus der Wirbelstruktur der WG folgen die Maxwell-Gleichungen der Gravitation:

\[
\begin{aligned}
\nabla \cdot \vec{E}_g &= -4\pi G \rho \\
\nabla \cdot \vec{B}_g &= 0 \\
\nabla \times \vec{E}_g &= -\frac{\partial \vec{B}_g}{\partial t} \\
\nabla \times \vec{B}_g &= \frac{4\pi G}{c^2} \vec{j} + \frac{1}{c^2} \frac{\partial \vec{E}_g}{\partial t}
\end{aligned}
\]

Dabei sind:
\begin{itemize}
\item \( \vec{E}_g = -\nabla \Phi \) – das gravito-elektrische Feld
\item \( \vec{B}_g = \nabla \times \vec{A}_g \) – das gravito-magnetische Feld
\item \( \vec{A}_g = \frac{2GM}{c^2 r} \vec{v} \) – das gravito-magnetische Vektorpotential
\end{itemize}

Die Maxwell-Gravitation ist eine lineare Theorie – sie enthält keine Raumzeitkrümmung.

\section{Gravitationswellen mit Dispersion}

Aus der Maxwell-Gravitation folgen Gravitationswellen.

Aber sie unterscheiden sich von denen der ART:

\[
\omega^2 = c^2 k^2 \cdot \left( \frac{k}{k_0} \right)^{2(D/3-1)}
\]

Mit \( D \approx 2,704 \) ist \( D/3 \approx 0,901 \), also:

\[
\omega \propto k^{0,901}
\]

Die Phasengeschwindigkeit wird frequenzabhängig:

\[
v_{\text{phase}}(\omega) = \frac{\omega}{k} \propto \omega^{1-D/3} \approx \omega^{-0,099}
\]

Dies ist eine testbare Vorhersage der Theorie.

\section{Die ART als Grenzfall}

Die ART ist der Grenzfall der Omega-Theorie für \( D \to 3 \).

\begin{itemize}
\item Für \( D = 3 \) wird die Dispersion linear: \( \omega = ck \).
\item Für \( D = 3 \) wird die Rotation zur Krümmung.
\item Für \( D = 3 \) wird die Maxwell-Gravitation zur ART.
\end{itemize}

Die ART ist eine hinreichende, aber keine notwendige Beschreibung der Gravitation.

\section{Zusammenfassung}

\begin{itemize}
\item Die Gravitation ist nicht Krümmung, sondern Rotation.
\item Die Omega-Theorie beschreibt Gravitation als Rotationsfeld \( \vec{\Omega} \).
\item Aus der Wirbelstruktur der WG folgt die Maxwell-Gravitation.
\item Gravitationswellen haben eine frequenzabhängige Dispersion: \( \omega \propto k^{D/3} \).
\item Die ART ist der Grenzfall \( D \to 3 \).
\end{itemize}